% Ejercicios de formulación orgánica (se incluye desde organica_apuntes_nuevo.tex).
% Cada ejercicio lleva su solución en un entorno solu: se escribe en ans1.tex
% y se imprime al final del documento.
% En las soluciones, \recomendado{...} indica el nombre según las recomendaciones
% IUPAC 2013 (Guía Breve 2021) cuando no coincide con el nombre tradicional.

\subsection*{Alcanos}

\begin{exercise}
    Formula los siguientes alcanos:
    \begin{multicols}{2}
        \begin{enumerate}[label=\alph*)]
            \item Pentano.
            \item 2,3,4-trimetilheptano.
            \item 4-etil-2,6-dimetiloctano.
            \item 4,6-dietil-2,4,8-trimetilnonano.
            \item 4-etil-2,2,5,8-tetrametil-6-propildecano.
            \item 3,7-dietil-5-isopropildecano.
        \end{enumerate}
    \end{multicols}
    \begin{solu}
        \textcolor{blue}{Se dibuja la cadena principal, se numera y se colocan los sustituyentes en sus posiciones.}
        \begin{enumerate}[label=\alph*)]
            \item \chemfig{CH_3-CH_2-CH_2-CH_2-CH_3} \quad (ya no se usa «n-pentano»)
            \item \chemfig{CH_3-CH(-[2]CH_3)-CH(-[6]CH_3)-CH(-[2]CH_3)-CH_2-CH_2-CH_3}
            \item \chemfig{CH_3-CH(-[2]CH_3)-CH_2-CH(-[6]CH_2-[0]CH_3)-CH_2-CH(-[2]CH_3)-CH_2-CH_3}
            \item \chemfig{CH_3-CH(-[2]CH_3)-CH_2-C(-[2]CH_3)(-[6]CH_2-[0]CH_3)-CH_2-CH(-[6]CH_2-[0]CH_3)-CH_2-CH(-[2]CH_3)-CH_3}
            \item \chemfig{CH_3-C(-[2]CH_3)(-[6]CH_3)-CH_2-CH(-[2]CH_2-[0]CH_3)-CH(-[6]CH_3)-CH(-[2]CH_2-[0]CH_2-[0]CH_3)-CH_2-CH(-[6]CH_3)-CH_2-CH_3}
            \item \chemfig{CH_3-CH_2-CH(-[2]CH_2-[0]CH_3)-CH_2-CH(-[6]CH(-[4]CH_3)-[6]CH_3)-CH_2-CH(-[2]CH_2-[0]CH_3)-CH_2-CH_2-CH_3}\\
                  \recomendado{3,7-dietil-5-(propan-2-il)decano}
        \end{enumerate}
    \end{solu}
\end{exercise}

\begin{exercise}
    Nombra los siguientes alcanos:
    \begin{enumerate}[label=\alph*)]
        \item \chemfig{CH_3-CH(-[2]CH_3)-CH_2-CH(-[6]CH_3)-CH(-[2]CH_3)-CH_3}
        \item \chemfig{CH_3-CH(-[6]CH_3)-CH_2-C(-[2]CH_3)(-[6]CH_3)-CH_2-CH_2-CH_3}
        \item \chemfig{CH_3-CH(-[2]CH_3)-CH(-[6]CH_3)-CH_2-C(-[2]CH_3)(-[6]CH_3)-CH_3}
        \item \chemfig{CH_3-C(-[2]CH_3)(-[6]CH_3)-CH_2-CH(-[6]CH_2-[0]CH_3)-CH(-[2]CH_3)-CH_2-CH(-[6]CH_3)-CH_3}
        \item \chemfig{CH_3-CH_2-CH(-[2]CH_3)-CH(-[6]CH_2-[6]CH_2-[6]CH_3)-CH(-[2]CH_3)-CH(-[6]CH_2-[0]CH_3)-CH(-[2]CH_3)-CH_3}
        \item \chemfig{CH_3-CH_2-C(-[2]CH_3)(-[6]CH(-[4]CH_3)-[6]CH_3)-CH(-[6]CH_3)-CH(-[6]CH_2-[0]CH_3)-CH_2-CH_2-CH_3}
    \end{enumerate}
    \begin{solu}
        \textcolor{blue}{Cadena más larga; si hay varias igual de largas, la de más sustituyentes. Después, localizadores más bajos para el conjunto de sustituyentes.}
        \begin{enumerate}[label=\alph*)]
            \item \textcolor{red}{2,3,5-trimetilhexano} (\{2,3,5\} es más bajo que \{2,4,5\})
            \item \textcolor{red}{2,4,4-trimetilheptano}
            \item \textcolor{red}{2,2,4,5-tetrametilhexano} (se numera desde la derecha: \{2,2,4,5\} es más bajo que \{2,3,5,5\})
            \item \textcolor{red}{4-etil-2,2,5,7-tetrametiloctano}
            \item \textcolor{red}{3-etil-2,4,6-trimetil-5-propiloctano}: hay varias cadenas de 8 carbonos y se elige la que tiene 5 sustituyentes; se numera desde la derecha.
            \item \textcolor{red}{3,5-dietil-2,3,4-trimetiloctano}: la cadena de 8 C que pasa por el grupo isopropilo tiene 5 sustituyentes (la horizontal solo tiene 4).
        \end{enumerate}
    \end{solu}
\end{exercise}

\subsection*{Alquenos y alquinos}

\begin{exercise}
    Formula los siguientes alquenos y alquinos:
    \begin{multicols}{2}
        \begin{enumerate}[label=\alph*)]
            \item But-1-eno.
            \item Pent-2-eno.
            \item Hexa-2,4-dieno.
            \item 3-butilhexa-1,4-dieno.
            \item But-2-ino.
            \item 3,4-dimetilpent-1-ino.
            \item 3,6-dimetilnona-1,4,7-triino.
            \item Pent-1-en-3-ino.
            \item Hept-3-en-1,6-diino.
            \item 4-etilhexa-1,3-dien-5-ino.
        \end{enumerate}
    \end{multicols}
    \begin{solu}
        \begin{enumerate}[label=\alph*)]
            \item \chemfig{CH_2=CH-CH_2-CH_3}
            \item \chemfig{CH_3-CH=CH-CH_2-CH_3}
            \item \chemfig{CH_3-CH=CH-CH=CH-CH_3}
            \item \chemfig{CH_2=CH-CH(-[6]CH_2-[0]CH_2-[0]CH_2-[0]CH_3)-CH=CH-CH_3}\\
                  \recomendado{4-eteniloct-2-eno}: la cadena más larga (8 C) pasa por el grupo butilo, aunque tenga un doble enlace menos.
            \item \chemfig{CH_3-C~C-CH_3}
            \item \chemfig{CH~C-CH(-[2]CH_3)-CH(-[6]CH_3)-CH_3}
            \item \chemfig{CH~C-CH(-[2]CH_3)-C~C-CH(-[2]CH_3)-C~C-CH_3}
            \item \chemfig{CH_2=CH-C~C-CH_3}
            \item \chemfig{CH~C-CH=CH-CH_2-C~CH} \quad (\{1,3,6\} es más bajo que \{1,4,6\})
            \item \chemfig{CH_2=CH-CH=C(-[6]CH_2-[0]CH_3)-C~CH} \quad (empate \{1,3,5\}: los dobles enlaces reciben 1 y 3)
        \end{enumerate}
    \end{solu}
\end{exercise}

\begin{exercise}
    Nombra los siguientes alquenos y alquinos:
    \begin{enumerate}[label=\alph*)]
        \item \chemfig{CH_3-CH=CH-CH_2-CH_2-CH_3}
        \item \chemfig{CH_2=CH-C(-[6]CH_3)=CH-CH_3}
        \item \chemfig{CH~C-CH(-[6]CH_3)-CH_2-CH_3}
        \item \chemfig{CH_3-CH=C(-[6]CH_3)-CH_2-CH_3}
        \item \chemfig{CH~C-C~C-CH_2-CH_3}
        \item \chemfig{CH~C-CH=CH-CH_3}
        \item \chemfig{CH_2=CH-CH(-[6]CH_3)-CH_2-CH_2-C~CH}
        \item \chemfig{CH_3-CH=C(-[6]CH_3)-CH_2-C~C-CH(-[2]CH_3)-CH=CH_2}
        \item \chemfig{CH_2=CH-CH(-[6]CH(-[6]CH_3)-CH_3)-CH_2-C~C-CH(-[6]CH_2-[6]CH_3)-C~CH}
    \end{enumerate}
    \begin{solu}
        \textcolor{blue}{Localizadores más bajos para el conjunto de dobles y triples enlaces; si hay empate, para los dobles.}
        \begin{enumerate}[label=\alph*)]
            \item \textcolor{red}{hex-2-eno}
            \item \textcolor{red}{3-metilpenta-1,3-dieno}
            \item \textcolor{red}{3-metilpent-1-ino}
            \item \textcolor{red}{3-metilpent-2-eno}
            \item \textcolor{red}{hexa-1,3-diino}
            \item \textcolor{red}{pent-3-en-1-ino} (\{1,3\} es más bajo que \{2,4\})
            \item \textcolor{red}{3-metilhept-1-en-6-ino} (empate \{1,6\}: el doble enlace recibe el 1)
            \item \textcolor{red}{3,7-dimetilnona-1,7-dien-4-ino} (se numera desde la derecha: \{1,4,7\})
            \item \textcolor{red}{3-etil-7-isopropilnon-8-en-1,4-diino}, \recomendado{3-etil-7-(propan-2-il)non-8-en-1,4-diino}. Todas las cadenas posibles tienen 9 C; se elige la que tiene más enlaces múltiples y se numera desde la derecha (\{1,4,8\}).
        \end{enumerate}
    \end{solu}
\end{exercise}

\subsection*{Hidrocarburos cíclicos}

\begin{exercise}
    Formula los siguientes hidrocarburos cíclicos:
    \begin{multicols}{2}
        \begin{enumerate}[label=\alph*)]
            \item Etilciclohexano.
            \item Ciclopenteno.
            \item Ciclohexino.
            \item 1,1,4,4-tetrametilciclohexano.
            \item 3-etilciclopenteno.
            \item 1,6-dimetilciclohexeno.
            \item 4-ciclobutilpent-1-ino.
            \item 3-ciclohexil-5-metilhex-2-eno.
            \item Ciclohexa-1,3-dieno.
            \item 3-ciclopentilprop-1-eno.
        \end{enumerate}
    \end{multicols}
    \begin{solu}
        \begin{enumerate}[label=\alph*)]
            \item \chemfig{*6(---(-CH_2-[0]CH_3)---)}
            \item \chemfig{*5(=----)}
            \item \chemfig{*6(~-----)}
            \item \chemfig{*6(-(-[::30]CH_3)(-[::-30]CH_3)---(-[::30]CH_3)(-[::-30]CH_3)--)}
            \item \chemfig{*5(=-(-CH_2-CH_3)---)}
            \item \chemfig{*6(-=(-CH_3)-(-CH_3)---)}\\
                  Los carbonos del doble enlace son el 1 y el 2; \{1,6\} es más bajo que \{2,3\}, por eso no se dice \mal{2,3-dimetilciclohexeno}.
            \item \chemfig{CH~C-CH_2-CH(-[6]*4(----))-CH_3}\\
                  \recomendado{(pent-4-in-2-il)ciclobutano}: el anillo tiene preferencia sobre la cadena.
            \item \chemfig{CH_3-CH=C(-[6]*6(------))-CH_2-CH(-[2]CH_3)-CH_3}\\
                  \recomendado{(5-metilhex-2-en-3-il)ciclohexano}
            \item \chemfig{*6(=-=---)}
            \item \chemfig{CH_2=CH-CH_2-[0]*5(-----)}\\
                  \recomendado{(prop-2-en-1-il)ciclopentano} o alilciclopentano
        \end{enumerate}
    \end{solu}
\end{exercise}

\begin{exercise}
    Nombra los siguientes hidrocarburos:
    \begin{enumerate}[label=\alph*)]
        \item \chemfig{*4(----)}
        \item \chemfig{*6(-----=)}
        \item \chemfig{*6(---(-CH_3)--(-CH_2-[2]CH_3)-)}
        \item \chemfig{*6(---=-=)}
        \item \chemfig{*6(---(-CH(-[:50]CH_3)(-[:-50]CH_3))---)}
        \item \chemfig{CH_3-CH(-[6]*6(------))-CH_2-CH_2-C(-[2]CH_3)(-[6]CH_3)-CH_2-CH_3}
        \item \chemfig{CH_3-CH(-[6]*5(-----))-CH=CH_2}
    \end{enumerate}
    \begin{solu}
        \begin{enumerate}[label=\alph*)]
            \item \textcolor{red}{ciclobutano}
            \item \textcolor{red}{ciclohexeno}
            \item \textcolor{red}{1-etil-3-metilciclohexano} (empate \{1,3\}: etil va antes por orden alfabético)
            \item \textcolor{red}{ciclohexa-1,3-dieno}
            \item \textcolor{red}{isopropilciclohexano}, \recomendado{(propan-2-il)ciclohexano}
            \item \textcolor{red}{2-ciclohexil-5,5-dimetilheptano}, \recomendado{(5,5-dimetilheptan-2-il)ciclohexano}
            \item \textcolor{red}{3-ciclopentilbut-1-eno}, \recomendado{(but-3-en-2-il)ciclopentano}
        \end{enumerate}
    \end{solu}
\end{exercise}

\subsection*{Hidrocarburos aromáticos}

\begin{exercise}
    Formula los siguientes hidrocarburos aromáticos:
    \begin{multicols}{2}
        \begin{enumerate}[label=\alph*)]
            \item Metilbenceno (tolueno).
            \item 1,3-dietilbenceno.
            \item Etenilbenceno (estireno).
            \item 1-butil-4-isopropilbenceno.
            \item \textit{p}-propiltolueno.
            \item 3-fenil-5-metilheptano.
            \item 4-fenilpent-1-eno.
            \item 2,4-difenil-3-metilhexano.
        \end{enumerate}
    \end{multicols}
    \begin{solu}
        \begin{enumerate}[label=\alph*)]
            \item \chemfig{CH_3-**6(------)}
            \item \chemfig{**6(-(-CH_2-CH_3)--(-CH_2-CH_3)---)}
            \item \chemfig{CH_2=CH-**6(------)}
            \item \chemfig{CH_3-CH_2-CH_2-CH_2-**6(---(-CH(-[2]CH_3)-CH_3)---)}\\
                  \recomendado{1-butil-4-(propan-2-il)benceno}
            \item \chemfig{CH_3-**6(---(-CH_2-CH_2-CH_3)---)}: \textcolor{red}{1-metil-4-propilbenceno}
            \item \chemfig{CH_3-CH_2-CH(-[6]**6(------))-CH_2-CH(-[2]CH_3)-CH_2-CH_3}\\
                  \recomendado{(5-metilheptan-3-il)benceno}
            \item \chemfig{CH_2=CH-CH_2-CH(-[6]**6(------))-CH_3}\\
                  \recomendado{(pent-4-en-2-il)benceno}
            \item \chemfig{CH_3-CH(-[6]**6(------))-CH(-[2]CH_3)-CH(-[6]**6(------))-CH_2-CH_3}\\
                  \recomendado{1,1\('\)-(3-metilhexano-2,4-diil)dibenceno}
        \end{enumerate}
    \end{solu}
\end{exercise}

\begin{exercise}
    Nombra los siguientes hidrocarburos aromáticos:
    \begin{enumerate}[label=\alph*)]
        \item \chemfig{CH_3-CH(-[6]**6(------))-CH_2-C(-[2]CH_3)(-[6]CH_3)-CH_3}
        \item \chemfig{**6(----(-CH_2-[:0]CH_3)--)}
        \item \chemfig{**6(-(-CH_2-CH_2-CH_3)--(-CH_3)---)}
        \item \chemfig{**6(-(-CH_3)--(-CH_3)--(-CH_3)-)}
        \item \chemfig{**6(----(-CH(-[1]CH_3)(-[3]H_3C))--)}
        \item \chemfig{**6(--(-[:0]CH_2-[:0]CH_2-[:0]CH_2-[:0]**6(------))----)}
        \item \chemfig{CH_3-CH=C(-[6]**6(------))-CH_2-CH_3}
    \end{enumerate}
    \begin{solu}
        \begin{enumerate}[label=\alph*)]
            \item \textcolor{red}{2,2-dimetil-4-fenilpentano}, \recomendado{(4,4-dimetilpentan-2-il)benceno}
            \item \textcolor{red}{etilbenceno}
            \item \textcolor{red}{1-metil-3-propilbenceno} (\textit{m}-propiltolueno)
            \item \textcolor{red}{1,3,5-trimetilbenceno}
            \item \textcolor{red}{isopropilbenceno} (cumeno), \recomendado{(propan-2-il)benceno}
            \item \textcolor{red}{1,3-difenilpropano}, \recomendado{1,1\('\)-(propano-1,3-diil)dibenceno}
            \item \textcolor{red}{3-fenilpent-2-eno}, \recomendado{(pent-2-en-3-il)benceno}
        \end{enumerate}
    \end{solu}
\end{exercise}

\subsection*{Derivados halogenados}

\begin{exercise}
    Formula los siguientes derivados halogenados:
    \begin{multicols}{2}
        \begin{enumerate}[label=\alph*)]
            \item 2-cloropropano.
            \item 1,3-dibromobenceno.
            \item 1,1,2,2-tetrafluoroetano.
            \item 1,4-diclorociclohexano.
            \item 4-bromopent-1-ino.
            \item 3-fluoro-5-metilhex-2-eno.
            \item 1,4-dibromo-6-ciclopentiloct-2-eno.
            \item 2,5-difenil-4-yodopent-1-eno.
            \item 4-clorobut-1-eno.
            \item 1,2-dibromobenceno.
        \end{enumerate}
    \end{multicols}
    \begin{solu}
        \begin{enumerate}[label=\alph*)]
            \item \chemfig{CH_3-CH(-[2]Cl)-CH_3}
            \item \chemfig{**6(-(-Br)--(-Br)---)}
            \item \chemfig{CHF_2-CHF_2}
            \item \chemfig{*6(-(-Cl)---(-Cl)--)}
            \item \chemfig{CH~C-CH_2-CH(-[2]Br)-CH_3}
            \item \chemfig{CH_3-CH=C(-[2]F)-CH_2-CH(-[2]CH_3)-CH_3}
            \item \chemfig{CH_2(-[2]Br)-CH=CH-CH(-[2]Br)-CH_2-CH(-[6]*5(-----))-CH_2-CH_3}\\
                  \recomendado{(5,8-dibromooct-6-en-3-il)ciclopentano}
            \item \chemfig{CH_2=C(-[2]**6(------))-CH_2-CH(-[6]I)-CH_2-[0]**6(------)}
            \item \chemfig{CH_2=CH-CH_2-CH_2-Cl}
            \item \chemfig{**6(-(-Br)-(-Br)----)}
        \end{enumerate}
    \end{solu}
\end{exercise}

\begin{exercise}
    Nombra los siguientes derivados halogenados:
    \begin{enumerate}[label=\alph*)]
        \item \chemfig{CH_3-I}
        \item \chemfig{CH_2(-[2]Cl)-CH_2(-[6]Cl)}
        \item \chemfig{CH_3-C(-[2]Br)(-[6]Br)-CH_2-CH_3}
        \item \chemfig{CH_2(-[2]Cl)-CH_2-CH=CH_2}
        \item \chemfig{CH~C-CH(-[2]Br)-CH_2-CH_2-Br}
        \item \chemfig{**6(----(-F)--)}
        \item \chemfig{CH_3-CH(-[2]Cl)-CH_2-CH_2-[0]*4(----)}
        \item \chemfig{*6(-(-Br)--(-Br)---)}
        \item \chemfig{CH_2(-[4]F)-CH(-[2]F)(-[6]F)}
        \item \chemfig{*5(-(-Cl)---(-Cl)-)}
    \end{enumerate}
    \begin{solu}
        \begin{enumerate}[label=\alph*)]
            \item \textcolor{red}{yodometano}
            \item \textcolor{red}{1,2-dicloroetano}
            \item \textcolor{red}{2,2-dibromobutano}
            \item \textcolor{red}{4-clorobut-1-eno} (el doble enlace se numera antes que el halógeno)
            \item \textcolor{red}{3,5-dibromopent-1-ino}
            \item \textcolor{red}{fluorobenceno}
            \item \textcolor{red}{1-ciclobutil-3-clorobutano}, \recomendado{(3-clorobutil)ciclobutano}
            \item \textcolor{red}{1,3-dibromociclohexano}
            \item \textcolor{red}{1,1,2-trifluoroetano}
            \item \textcolor{red}{1,3-diclorociclopentano}
        \end{enumerate}
    \end{solu}
\end{exercise}

\subsection*{Alcoholes y éteres}

\begin{exercise}
    Formula los siguientes alcoholes y éteres:
    \begin{multicols}{2}
        \begin{enumerate}[label=\alph*)]
            \item 3-metilpentan-1-ol.
            \item Butano-1,2,3-triol.
            \item 2-fenilpropano-1,3-diol.
            \item Ciclohexanol.
            \item Hexa-3,5-dien-2-ol.
            \item Fenol (bencenol).
            \item 2-etilpentan-1-ol.
            \item Pent-3-en-1-ol.
            \item Etil isopropil éter.
            \item Etenil fenil éter.
            \item Dimetil éter.
            \item Butil ciclopentil éter.
        \end{enumerate}
    \end{multicols}
    \begin{solu}
        \begin{enumerate}[label=\alph*)]
            \item \chemfig{CH_2(-[2]OH)-CH_2-CH(-[2]CH_3)-CH_2-CH_3}
            \item \chemfig{CH_2(-[2]OH)-CH(-[2]OH)-CH(-[2]OH)-CH_3}
            \item \chemfig{CH_2(-[2]OH)-CH(-[6]**6(------))-CH_2(-[2]OH)}
            \item \chemfig{*6(---(-OH)---)}
            \item \chemfig{CH_3-CH(-[2]OH)-CH=CH-CH=CH_2}
            \item \chemfig{**6(---(-OH)---)}
            \item \chemfig{CH_2(-[2]OH)-CH(-[6]CH_2-[0]CH_3)-CH_2-CH_2-CH_3}
            \item \chemfig{CH_2(-[2]OH)-CH_2-CH=CH-CH_3}
            \item \chemfig{CH_3-CH_2-O-CH(-[2]CH_3)-CH_3}: \recomendado{2-etoxipropano}
            \item \chemfig{CH_2=CH-O-**6(------)}: \recomendado{(eteniloxi)benceno}
            \item \chemfig{CH_3-O-CH_3}: \recomendado{metoximetano}
            \item \chemfig{CH_3-CH_2-CH_2-CH_2-O-*5(-----)}: \recomendado{butoxiciclopentano} (el anillo es el progenitor)
        \end{enumerate}
    \end{solu}
\end{exercise}

\begin{exercise}
    Nombra los siguientes alcoholes y éteres:
    \begin{enumerate}[label=\alph*)]
        \item \chemfig{CH_3-OH}
        \item \chemfig{CH_2(-[2]OH)-CH_2-CH(-[6]CH_2-[0]CH_3)-CH=CH_2}
        \item \chemfig{CH_3-CH(-[2]OH)-CH(-[6]Br)-CH_2(-[2]OH)}
        \item \chemfig{CH_3-CH_2-CH(-[2]CH_3)-CH(-[6]CH_3)-CH(-[2]OH)-CH=CH_2}
        \item \chemfig{CH_2=CH-CH(-[2]OH)-CH_2(-[2]OH)}
        \item \chemfig{CH_3-CH_2-CH_2-O-CH_2-CH_2-CH_3}
        \item \chemfig{CH_2=CH-O-CH(-[2]CH_3)-CH_3}
        \item \chemfig{CH_3-CH_2-CH_2-CH_2-CH_2-O-C~CH}
    \end{enumerate}
    \begin{solu}
        \begin{enumerate}[label=\alph*)]
            \item \textcolor{red}{metanol}
            \item \textcolor{red}{3-etilpent-4-en-1-ol}: las dos cadenas que contienen el \ch{-OH} tienen 5 C; se elige la que tiene el doble enlace.
            \item \textcolor{red}{2-bromobutano-1,3-diol}
            \item \textcolor{red}{4,5-dimetilhept-1-en-3-ol}: el \ch{-OH} recibe el localizador más bajo (3).
            \item \textcolor{red}{but-3-eno-1,2-diol}
            \item \textcolor{red}{dipropil éter}, \recomendado{1-propoxipropano}
            \item \textcolor{red}{etenil isopropil éter}, \recomendado{2-(eteniloxi)propano}: el propano (3 C) es más largo que el eteno.
            \item \textcolor{red}{etinil pentil éter}, \recomendado{1-(etiniloxi)pentano}
        \end{enumerate}
    \end{solu}
\end{exercise}

\subsection*{Aldehídos y cetonas}

\begin{exercise}
    Formula los siguientes aldehídos y cetonas:
    \begin{multicols}{2}
        \begin{enumerate}[label=\alph*)]
            \item Etanal (acetaldehído).
            \item 3-metilpentanal.
            \item Propenal.
            \item 5-ciclohexilpent-3-inal.
            \item Hex-2-enodial.
            \item Hexano-2,4-diona.
            \item 1,4-difenilpentan-2-ona.
            \item Benzaldehído.
            \item 2-metilpentanodial.
            \item 3-metilpent-2-enal.
            \item Pentan-2-ona.
            \item 3-clorobutan-2-ona.
            \item Hexa-1,5-dien-3-ona.
        \end{enumerate}
    \end{multicols}
    \begin{solu}
        \begin{enumerate}[label=\alph*)]
            \item \chemfig{CH_3-!{aldehido}}
            \item \chemfig{CH_3-CH_2-CH(-[2]CH_3)-CH_2-!{aldehido}}
            \item \chemfig{CH_2=CH-!{aldehido}} (prop-2-enal)
            \item \chemfig{!{aldehido2}-CH_2-C~C-CH_2-[0]*6(------)}
            \item \chemfig{!{aldehido2}-CH=CH-CH_2-CH_2-!{aldehido}}
            \item \chemfig{CH_3-C(=[2]O)-CH_2-C(=[2]O)-CH_2-CH_3} (no «hexa-»: no tiene dobles enlaces C=C)
            \item \chemfig{CH_3-CH(-[6]**6(------))-CH_2-C(=[2]O)-CH_2-[0]**6(------)}
            \item \chemfig{**6(---(-CHO)---)}
            \item \chemfig{!{aldehido2}-CH(-[6]CH_3)-CH_2-CH_2-!{aldehido}}
            \item \chemfig{CH_3-CH_2-C(-[6]CH_3)=CH-!{aldehido}}
            \item \chemfig{CH_3-C(=[2]O)-CH_2-CH_2-CH_3}
            \item \chemfig{CH_3-C(=[2]O)-CH(-[2]Cl)-CH_3}
            \item \chemfig{CH_2=CH-C(=[2]O)-CH_2-CH=CH_2}
        \end{enumerate}
    \end{solu}
\end{exercise}

\begin{exercise}
    Nombra los siguientes aldehídos y cetonas:
    \begin{enumerate}[label=\alph*)]
        \item \chemfig{H-!{aldehido}}
        \item \chemfig{CH_3-CH_2-CH_2-!{aldehido}}
        \item \chemfig{!{aldehido2}-CH=CH-!{aldehido}}
        \item \chemfig{CH_2=C(-[6]!{fenil})-CH_2-CH_2-!{aldehido}}
        \item \chemfig{!{aldehido2}-CH=CH-CH_2-CH(-[2]CH_3)-!{aldehido}}
        \item \chemfig{CH_3-CH(-[6]!{fenil})-CH=CH-!{aldehido}}
        \item \chemfig{!{aldehido2}-CH_2-C~C-CH_2-CH_2-!{aldehido}}
        \item \chemfig{CH_3-C(=[2]O)-CH_2-CH_3}
        \item \chemfig{CH_3-CH=CH-CH_2-C(=[2]O)-CH_3}
        \item \chemfig{CH_3-C(=[2]O)-CH_2-CH_2-CH_2-C(=[2]O)-CH_3}
        \item \chemfig{CH_3-CH(-[2]CH_3)-C(=[6]O)-CH_2-CH(-[2]CH_3)-CH_3}
        \item \chemfig{CH_2=CH-C(=[2]O)-CH=CH-CH_3}
    \end{enumerate}
    \begin{solu}
        \begin{enumerate}[label=\alph*)]
            \item \textcolor{red}{metanal} (formaldehído)
            \item \textcolor{red}{butanal}
            \item \textcolor{red}{but-2-enodial}
            \item \textcolor{red}{4-fenilpent-4-enal}
            \item \textcolor{red}{5-metilhex-2-enodial}: los dos \ch{-CHO} dan 1 y 6 en ambos sentidos; decide el doble enlace (2 mejor que 4).
            \item \textcolor{red}{4-fenilpent-2-enal}
            \item \textcolor{red}{hept-3-inodial}
            \item \textcolor{red}{butanona} (butan-2-ona, etil metil cetona)
            \item \textcolor{red}{hex-4-en-2-ona}
            \item \textcolor{red}{heptano-2,6-diona}
            \item \textcolor{red}{2,5-dimetilhexan-3-ona}
            \item \textcolor{red}{hexa-1,4-dien-3-ona}
        \end{enumerate}
    \end{solu}
\end{exercise}

\subsection*{Ácidos carboxílicos y ésteres}

\begin{exercise}
    Formula los siguientes ácidos y ésteres:
    \begin{multicols}{2}
        \begin{enumerate}[label=\alph*)]
            \item Ácido etanoico (ácido acético).
            \item Ácido 3-metilhexanoico.
            \item Ácido 2-fenilpentanodioico.
            \item Ácido tricloroetanoico.
            \item Ácido but-3-enoico.
            \item Ácido hepta-2,4-dienoico.
            \item Ácido pent-2-enodioico.
            \item Ácido benzoico.
            \item Butanoato de metilo.
            \item Propanoato de etilo.
            \item Benzoato de propilo.
            \item Etanoato de octilo.
            \item 3-cloropentanoato de etenilo.
            \item But-3-enoato de isopropilo.
        \end{enumerate}
    \end{multicols}
    \begin{solu}
        \begin{enumerate}[label=\alph*)]
            \item \chemfig{CH_3-!{acido}}
            \item \chemfig{CH_3-CH_2-CH_2-CH(-[2]CH_3)-CH_2-!{acido}}
            \item \chemfig{!{acido2}-CH(-[6]**6(------))-CH_2-CH_2-!{acido}}
            \item \chemfig{CCl_3-!{acido}}
            \item \chemfig{CH_2=CH-CH_2-!{acido}}
            \item \chemfig{CH_3-CH_2-CH=CH-CH=CH-!{acido}}
            \item \chemfig{!{acido2}-CH=CH-CH_2-!{acido}}
            \item \chemfig{**6(---(-COOH)---)}
            \item \chemfig{CH_3-CH_2-CH_2-C(=[2]O)-O-CH_3}
            \item \chemfig{CH_3-CH_2-C(=[2]O)-O-CH_2-CH_3}
            \item \chemfig{CH_3-CH_2-CH_2-O-C(=[2]O)-**6(------)}
            \item \ch{CH3-COO-(CH2)7-CH3}
            \item \chemfig{CH_3-CH_2-CH(-[2]Cl)-CH_2-C(=[2]O)-O-CH=CH_2}
            \item \chemfig{CH_2=CH-CH_2-C(=[2]O)-O-CH(-[2]CH_3)-CH_3}\\
                  \recomendado{but-3-enoato de propan-2-ilo}
        \end{enumerate}
    \end{solu}
\end{exercise}

\begin{exercise}
    Nombra los siguientes ácidos y ésteres:
    \begin{enumerate}[label=\alph*)]
        \item \chemfig{CH_3-CH_2-CH_2-!{acido}}
        \item \chemfig{!{acido2}-CH_2-CH_2-!{acido}}
        \item \chemfig{CH_3-CH=CH-CH_2-!{acido}}
        \item \chemfig{CH_3-CH(-[2]CH_3)-CH_2-CH(-[2]CH_3)-CH_2-!{acido}}
        \item \chemfig{!{acido2}-C~C-CH_2-!{acido}}
        \item \chemfig{CH_3-CH_2-CH(-[6]!{ciclopentil})-!{acido}}
        \item \chemfig{CH_3-C(=[2]O)-O-CH_3}
        \item \chemfig{CH_3-CH_2-CH_2-C(=[2]O)-O-CH_2-CH_3}
        \item \chemfig{CH_3-CH_2-CH_2-CH_2-CH_2-CH_2-CH_2-C(=[2]O)-O-CH_2-CH_2-CH_3}
        \item \chemfig{CH_3-CH_2-CH_2-CH_2-CH_2-C(=[2]O)-O-CH_3}
        \item \chemfig{CH_3-CH(-[2]Br)-CH_2-C(=[2]O)-O-CH=CH_2}
        \item \chemfig{CH_2(-[:180]!{fenil})-CH_2-C(=[2]O)-O-CH_2-CH_2-CH_2-CH_2-CH_2-CH_3}
    \end{enumerate}
    \begin{solu}
        \textcolor{blue}{El carbono del grupo \ch{-COOH} (o \ch{-COO-}) es siempre el carbono 1.}
        \begin{enumerate}[label=\alph*)]
            \item \textcolor{red}{ácido butanoico}
            \item \textcolor{red}{ácido butanodioico} (ácido succínico)
            \item \textcolor{red}{ácido pent-3-enoico}
            \item \textcolor{red}{ácido 3,5-dimetilhexanoico}
            \item \textcolor{red}{ácido pent-2-inodioico}
            \item \textcolor{red}{ácido 2-ciclopentilbutanoico} (el \ch{-COOH} está en la cadena: la cadena es el progenitor)
            \item \textcolor{red}{etanoato de metilo} (acetato de metilo)
            \item \textcolor{red}{butanoato de etilo}
            \item \textcolor{red}{octanoato de propilo}
            \item \textcolor{red}{hexanoato de metilo}
            \item \textcolor{red}{3-bromobutanoato de etenilo} (de vinilo)
            \item \textcolor{red}{3-fenilpropanoato de hexilo}
        \end{enumerate}
    \end{solu}
\end{exercise}

\subsection*{Compuestos nitrogenados}

\begin{exercise}
    Formula los siguientes compuestos nitrogenados:
    \begin{multicols}{2}
        \begin{enumerate}[label=\alph*)]
            \item Isopropilamina.
            \item Pentan-3-amina.
            \item Butano-1,3-diamina.
            \item Etano-1,2-diamina.
            \item 3-etilhexan-3-amina.
            \item 3,5-dimetilhexan-1-amina.
            \item Pent-3-en-2-amina.
            \item \textit{N}-etil-\textit{N}-metilpropan-1-amina.
            \item \textit{N}-metilanilina.
            \item \textit{N}-ciclopentilbutan-1-amina.
            \item Etanamida.
            \item \textit{N}-metiletanamida.
            \item 4-fenilpentanamida.
            \item \textit{N}-etilhex-4-enamida.
        \end{enumerate}
    \end{multicols}
    \begin{solu}
        \begin{enumerate}[label=\alph*)]
            \item \chemfig{CH_3-CH(-[2]NH_2)-CH_3}: \recomendado{propan-2-amina}
            \item \chemfig{CH_3-CH_2-CH(-[2]NH_2)-CH_2-CH_3}
            \item \chemfig{CH_2(-[2]NH_2)-CH_2-CH(-[2]NH_2)-CH_3}
            \item \chemfig{H_2N-CH_2-CH_2-NH_2}
            \item \chemfig{CH_3-CH_2-C(-[2]NH_2)(-[6]CH_2-[0]CH_3)-CH_2-CH_2-CH_3}
            \item \chemfig{H_2N-CH_2-CH_2-CH(-[2]CH_3)-CH_2-CH(-[2]CH_3)-CH_3}
            \item \chemfig{CH_3-CH(-[2]NH_2)-CH=CH-CH_3}
            \item \chemfig{CH_3-CH_2-CH_2-N(-[2]CH_3)-CH_2-CH_3}
            \item \chemfig{**6(---(-NH-CH_3)---)}: \recomendado{\textit{N}-metilanilina} (\textit{N}-metilbencenamina)
            \item \chemfig{CH_3-CH_2-CH_2-CH_2-NH-[0]*5(-----)}: \recomendado{\textit{N}-butilciclopentanamina} (el anillo es el progenitor)
            \item \chemfig{CH_3-C(=[2]O)-NH_2}
            \item \chemfig{CH_3-C(=[2]O)-NH-CH_3}
            \item \chemfig{CH_3-CH(-[6]**6(------))-CH_2-CH_2-C(=[2]O)-NH_2}
            \item \chemfig{CH_3-CH=CH-CH_2-CH_2-C(=[2]O)-NH-CH_2-CH_3}
        \end{enumerate}
    \end{solu}
\end{exercise}

\begin{exercise}
    Nombra los siguientes compuestos nitrogenados:
    \begin{enumerate}[label=\alph*)]
        \item \chemfig{CH_3-CH(-[6]NH_2)-CH_2-CH_3}
        \item \chemfig{CH_3-CH_2-CH_2-NH_2}
        \item \chemfig{CH_3-CH(-[6]NH_2)-CH_2-CH(-[6]NH_2)-CH_2-CH_2-NH_2}
        \item \chemfig{CH_3-CH(-[6]CH_3)-NH-CH=CH_2}
        \item \chemfig{CH_3-CH_2-CH_2-CH_2-CH_2-C(=[2]O)-NH_2}
        \item \chemfig{CH_3-CH_2-C(=[2]O)-N(-[6]CH_3)-CH_3}
        \item \chemfig{CH_3-CH=CH-CH_2-C(=[2]O)-NH_2}
        \item \chemfig{CH_3-CH_2-CH(-[2]Br)-CH_2-CH_2-C(=[2]O)-NH-CH_3}
    \end{enumerate}
    \begin{solu}
        \begin{enumerate}[label=\alph*)]
            \item \textcolor{red}{butan-2-amina}
            \item \textcolor{red}{propan-1-amina} (propilamina)
            \item \textcolor{red}{hexano-1,3,5-triamina} (se numera desde la derecha: \{1,3,5\})
            \item \textcolor{red}{\textit{N}-etenilpropan-2-amina}: el propano (3 C) es más largo que el eteno.
            \item \textcolor{red}{hexanamida}
            \item \textcolor{red}{\textit{N},\textit{N}-dimetilpropanamida}
            \item \textcolor{red}{pent-3-enamida}
            \item \textcolor{red}{4-bromo-\textit{N}-metilhexanamida}
        \end{enumerate}
    \end{solu}
\end{exercise}

\subsection*{Nitrilos y nitroderivados}

\begin{exercise}
    Formula los siguientes compuestos:
    \begin{multicols}{2}
        \begin{enumerate}[label=\alph*)]
            \item Butanonitrilo.
            \item Benzonitrilo.
            \item 2-metilpropanonitrilo.
            \item Nitrobenceno.
            \item 2-nitropropano.
            \item 3-nitropent-1-eno.
        \end{enumerate}
    \end{multicols}
    \begin{solu}
        \begin{enumerate}[label=\alph*)]
            \item \chemfig{CH_3-CH_2-CH_2-C~N}
            \item \chemfig{**6(---(-C~N)---)}
            \item \chemfig{CH_3-CH(-[2]CH_3)-C~N}
            \item \chemfig{**6(---(-NO_2)---)}
            \item \chemfig{CH_3-CH(-[2]NO_2)-CH_3}
            \item \chemfig{CH_2=CH-CH(-[2]NO_2)-CH_2-CH_3}
        \end{enumerate}
    \end{solu}
\end{exercise}

\begin{exercise}
    Nombra los siguientes compuestos:
    \begin{enumerate}[label=\alph*)]
        \item \chemfig{CH_3-CH_2-C~N}
        \item \chemfig{N~C-CH_2-CH_2-C~N}
        \item \chemfig{CH_3-NO_2}
        \item \chemfig{CH_3-CH(-[2]CH_3)-CH_2-C~N}
        \item \chemfig{**6(-(-NO_2)--(-CH_3)---)}
    \end{enumerate}
    \begin{solu}
        \begin{enumerate}[label=\alph*)]
            \item \textcolor{red}{propanonitrilo}
            \item \textcolor{red}{butanodinitrilo}
            \item \textcolor{red}{nitrometano}
            \item \textcolor{red}{3-metilbutanonitrilo} (el C del \ch{-C+N} es el 1)
            \item \textcolor{red}{1-metil-3-nitrobenceno} (empate \{1,3\}: metil va antes por orden alfabético)
        \end{enumerate}
    \end{solu}
\end{exercise}

\subsection*{Compuestos con varios grupos funcionales}

\begin{exercise}
    Formula los siguientes compuestos:
    \begin{multicols}{2}
        \begin{enumerate}[label=\alph*)]
            \item Ácido 2-hidroxipropanoico (ácido láctico).
            \item 4-hidroxibutan-2-ona.
            \item 3-oxobutanal.
            \item Ácido 2-aminopropanoico (alanina).
            \item Ácido 4-oxopentanoico.
        \end{enumerate}
    \end{multicols}
    \begin{solu}
        \begin{enumerate}[label=\alph*)]
            \item \chemfig{CH_3-CH(-[2]OH)-!{acido}}
            \item \chemfig{HO-CH_2-CH_2-C(=[2]O)-CH_3}
            \item \chemfig{CH_3-C(=[2]O)-CH_2-!{aldehido}}
            \item \chemfig{CH_3-CH(-[2]NH_2)-!{acido}}
            \item \chemfig{CH_3-C(=[2]O)-CH_2-CH_2-!{acido}}
        \end{enumerate}
    \end{solu}
\end{exercise}

\begin{exercise}
    Nombra los siguientes compuestos:
    \begin{enumerate}[label=\alph*)]
        \item \chemfig{HO-CH_2-CH_2-!{aldehido}}
        \item \chemfig{H_2N-CH_2-CH_2-OH}
        \item \chemfig{CH_3-C(=[2]O)-CH_2-!{acido}}
        \item \chemfig{N~C-CH_2-!{acido}}
        \item \chemfig{CH_3-CH(-[2]OH)-CH_2-C(=[2]O)-CH_3}
    \end{enumerate}
    \begin{solu}
        \textcolor{blue}{El grupo de mayor jerarquía da el sufijo; el resto se nombran como prefijos.}
        \begin{enumerate}[label=\alph*)]
            \item \textcolor{red}{3-hidroxipropanal} (aldehído \(>\) alcohol)
            \item \textcolor{red}{2-aminoetan-1-ol} (alcohol \(>\) amina)
            \item \textcolor{red}{ácido 3-oxobutanoico} (ácido \(>\) cetona)
            \item \textcolor{red}{ácido cianoetanoico} (ácido cianoacético; ácido \(>\) nitrilo)
            \item \textcolor{red}{4-hidroxipentan-2-ona} (cetona \(>\) alcohol)
        \end{enumerate}
    \end{solu}
\end{exercise}
